\chapter{Introduction to De Novo Peptide Sequencing and Mass Spectrometry}

Proteins are the primary functional macromolecular effectors of biological processes, executing enzymatic catalysis, structural scaffolding, signal transduction, and immunological defense within all living organisms. Deciphering the exact amino acid sequence of proteins and peptides is therefore foundational to molecular biology, pharmacology, clinical oncology, and biotechnology. Tandem mass spectrometry (MS/MS) coupled with liquid chromatography (LC-MS/MS) has emerged as the premier high-throughput analytical technology for resolving complex cellular proteomes. However, translating high-dimensional, noisy fragmentation spectra into accurate primary sequences remains one of the most demanding inverse problems in computational biology.

\section{The Proteomics Frontier and Tandem Mass Spectrometry}

In a standard bottom-up proteomics workflow, extracted proteins are enzymatically digested into shorter peptide fragments (typically 7 to 35 amino acids in length) using endopeptidases such as trypsin \citep{Olsen2004, Vandermarliere2013}. These peptides are separated by reversed-phase high-performance liquid chromatography and introduced into a mass spectrometer via electrospray ionization (ESI). 

During the initial survey stage (MS1), the instrument measures the mass-to-charge ratios ($m/z$) and isotopic abundance envelopes of intact, multiply protonated peptide precursor ions. Subsequently, selected precursor ions within narrow isolation windows are accelerated into a collision cell for tandem mass spectrometry (MS2). In Higher-Energy Collisional Dissociation (HCD) \citep{Olsen2007}, precursor ions collide with neutral inert gas molecules (such as nitrogen or argon), converting kinetic energy into vibrational excitation. This vibrational excitation induces unimolecular thermal decomposition across the weakest covalent bonds of the peptide backbone, predominantly the amide peptide bonds \citep{Paizs2005}. The resulting fragment ions are injected into an ultra-high-resolution mass analyzer, such as an Orbitrap \citep{Makarov2000, Zubarev2007}, which records their exact $m/z$ and intensity values, generating an experimental MS2 fragmentation spectrum.

\section{Database Search Engines versus De Novo Sequencing}

For over two decades, the dominant paradigm for analyzing tandem mass spectra has been database-dependent search algorithms, pioneered by SEQUEST \citep{Eng1994} and subsequently refined in tools such as Mascot \citep{Perkins1999}, MaxQuant \citep{Cox2008}, and MSFragger \citep{Kong2017}, where spectrum-to-sequence matches are statistically validated using target-decoy false discovery rate (FDR) estimation \citep{Elias2007}. Database search engines operate by matching experimental spectra against \emph{in silico} theoretical spectra generated from reference proteome sequence databases (such as UniProt). While highly sensitive and computationally tractable for well-characterized model organisms, database searching possesses severe intrinsic limitations \citep{Nesvizhskii2014}:

\begin{enumerate}[(i)]
\item \textbf{Unannotated Genomes and Non-Model Species:} For the vast majority of biological species on Earth, high-quality reference genomes or proteome annotations do not exist. Database searching fails entirely when the target organism is unsequenced \citep{Nesvizhskii2014, Aebersold2016}.
\item \textbf{Cancer Neoantigens and Somatic Mutations:} Malignant cells generate novel peptide sequences resulting from point mutations, gene fusions, alternative splicing, and frame-shift events that are absent from canonical reference databases. Identifying these tumor-specific antigens is vital for personalized immunotherapy and cancer vaccine development \citep{Schumacher2015, BassaniSternberg2016}.
\item \textbf{Immunopeptidomics and Antibody Repertoires:} B-cell and T-cell receptor repertoires undergo somatic hypermutation and $V(D)J$ genetic recombination, yielding hypervariable antigen-binding complementarity-determining regions (CDRs). Assembling monoclonal antibodies or polyclonal serum repertoires directly from biological fluids requires sequence identification without genomic references \citep{Purcell2019, Bonissone2016}.
\item \textbf{Post-Translational Modifications (PTMs):} Variable chemical modifications (phosphorylation, acetylation, ubiquitination, methylation) expand the combinatorial search space exponentially, rendering brute-force database searches computationally intractable \citep{Chick2015, Kong2017}.
\end{enumerate}

To overcome these barriers, \emph{de novo} peptide sequencing aims to reconstruct the exact primary amino acid sequence directly from the experimental spectrum $\mathcal{S}$ and precursor mass $M_{\text{prec}}$, relying solely on the physical laws of mass conservation and fragmentation chemistry.

\section{Physical Principles of Collisional Fragmentation}

To understand the mathematical structure of the sequencing problem, we examine the nomenclature and mechanisms governing peptide fragmentation. Following the standardized Roepstorff-Biemann nomenclature \citep{Roepstorff1984, Biemann1988}, cleavage of the amide bond ($-CO-NH-$) along the peptide backbone generates two complementary series of fragment ions:

\begin{defn}[Prefix and Suffix Fragment Ions]
Let $Y = (y_1, y_2, \dots, y_L)$ denote a peptide of length $L$, where each $y_i \in \mathcal{V}_{\text{AA}}$ is an amino acid residue with monoisotopic mass $m(y_i)$. 
\begin{enumerate}[(i)]
\item A \textbf{$b$-ion} of length $k \in \{1, \dots, L-1\}$ contains the N-terminal prefix sequence $(y_1, \dots, y_k)$. Its theoretical monoisotopic mass is:
\begin{equation}
m(b_k) = \sum_{j=1}^k m(y_j) + m(\text{H}^+)
\label{eq:b_ion_mass}
\end{equation}
where $m(\text{H}^+) \approx 1.007276\,\text{Da}$ is the proton mass.
\item A \textbf{$y$-ion} of length $k \in \{1, \dots, L-1\}$ contains the C-terminal suffix sequence $(y_{L-k+1}, \dots, y_L)$. Its theoretical monoisotopic mass is:
\begin{equation}
m(y_k) = \sum_{j=L-k+1}^L m(y_j) + m(\text{H}_2\text{O}) + m(\text{H}^+)
\label{eq:y_ion_mass}
\end{equation}
where $m(\text{H}_2\text{O}) \approx 18.010565\,\text{Da}$ represents the terminal water molecule retained by the carboxylic acid terminus.
\end{enumerate}
\end{defn}

Because cleavage of a specific peptide bond produces either a $b$-ion or a $y$-ion, every $b_k$ ion has a complementary $y_{L-k}$ ion satisfying the fundamental mass conservation duality:
\begin{equation}
m(b_k) + m(y_{L-k}) = M_{\text{prec}} + m(\text{H}^+)
\label{eq:complementary_ion_duality}
\end{equation}
where $M_{\text{prec}} = \sum_{i=1}^L m(y_i) + m(\text{H}_2\text{O}) + m(\text{H}^+)$ is the singly protonated intact precursor mass. The mass difference between two consecutive fragment ions of the same series reveals the mass of the intervening amino acid:
\begin{equation}
m(b_{k}) - m(b_{k-1}) = m(y_{L-k+1}) - m(y_{L-k}) = m(y_k)
\label{eq:mass_ladder}
\end{equation}
In an idealized, noise-free mass spectrum, a complete ladder of $b$-ions or $y$-ions would permit trivial sequence extraction by reading off adjacent mass intervals. However, real experimental spectra deviate substantially from this idealization due to missing fragmentation peaks, internal fragments, neutral losses ($\text{H}_2\text{O}$, $\text{NH}_3$), multiply charged ions, chemical noise, and instrument measurement jitter \citep{Paizs2005}.

\section{Computational Complexity and Predecessor Algorithms}

The primary amino acid alphabet consists of 20 canonical residues with monoisotopic masses ranging from Glycine ($57.021\,\text{Da}$) to Tryptophan ($186.079\,\text{Da}$). For a typical tryptic peptide of length $L \approx 15$, the number of possible sequence permutations is $20^{15} \approx 3.28 \times 10^{19}$. Filtering by precursor neutral mass restricts the candidate space, but still leaves millions of isobaric sequence permutations.

Early computational approaches framed \emph{de novo} sequencing as finding the highest-scoring path through a directed acyclic spectrum graph \citep{Dancik1999, Chen2001, Frank2005}. Nodes represented observed experimental peaks, and directed edges corresponded to valid amino acid residue masses. While dynamic programming could compute optimal paths in polynomial time, these heuristic graph methods struggled with missing peaks, false ladder connections, and non-linear peak intensity variations.

Recent advances in deep learning have significantly improved \emph{de novo} sequencing performance. DeepNovo \citep{Tran2017} combined convolutional neural networks (CNNs) for peak feature extraction with long short-term memory (LSTM) networks and beam search decoding. PointNovo \citep{Qiao2021} introduced continuous order-invariant spectrum encodings. Casanovo \citep{Yilmaz2022} adapted the Transformer architecture, framing sequencing as an autoregressive machine translation problem from peak coordinates to amino acid tokens. More recently, InstaNovo \citep{Eloff2025} incorporated mass-guided beam search with deep autoregressive Transformers, setting a new benchmark for sequencing accuracy across large-scale human proteomes.

Despite their success, existing autoregressive models suffer from fundamental algorithmic limitations:
\begin{enumerate}[(i)]
\item \textbf{Sequential Decoding Latency:} Generating a sequence of length $L$ requires $L$ sequential forward passes through the Transformer decoder. With typical peptide lengths of $L \in [7, 30]$ and beam sizes of $B \in [5, 10]$, decoding speed is capped at 25 to 55 spectra per second on modern GPUs, creating an acute bottleneck for clinical high-throughput repositories containing tens of millions of spectra.
\item \textbf{Error Propagation and Exposure Bias:} Because autoregressive models generate tokens left-to-right, an error made on an early residue propagates through the entire remaining sequence, leading to sequence-level failure.
\item \textbf{Directional Ion Asymmetry:} Left-to-right generation favors $b$-ions over $y$-ions, violating the physical bilateral symmetry of MS/MS fragmentation expressed in Equation~\ref{eq:complementary_ion_duality}.
\item \textbf{Discretization Failure in Continuous Flows:} Continuous normalizing flow approaches, such as PowerNovo2 \citep{Petrovskiy2026}, model continuous latent vectors rather than discrete sequences. Projecting continuous vectors onto discrete amino acid masses causes massive discretization error, resulting in low exact match rates (3.16\% on standard biological benchmarks).
\end{enumerate}

\section{Outline of Contributions}

To simultaneously overcome the latency bottleneck of autoregressive models and the discretization failure of continuous normalizing flows, this thesis introduces \textbf{DFlowNovo}, a continuous-time Discrete Flow Matching framework tailored for high-throughput \emph{de novo} peptide sequencing. 

The primary contributions of this thesis are:
\begin{enumerate}[(1)]
\item \textbf{Discrete Flow Formulation for Proteomics:} We establish the first continuous-time Discrete Flow Matching (DFM) formulation directly defined over the discrete probability simplex of amino acid residues, avoiding continuous discretization errors entirely.
\item \textbf{Exact Polynomial-Time Knapsack Guidance (\texttt{KnapsackDP}):} We derive and implement a GPU-accelerated dynamic programming reachability filter that prunes mathematically invalid amino acid transitions during flow integration in $\mathcal{O}(1)$ time per token, enforcing 100\% mass compliance.
\item \textbf{Empirical State-of-the-Art Accuracy:} Across the standardized 104k Nine-Species biological benchmark \citep{Tran2017} (published on the Hugging Face Hub as \texttt{InstaDeepAI/ms\_ninespecies\_benchmark}), DFlowNovo achieves 68.28\% strict exact match and 83.01\% residue F1 score, outperforming InstaNovo v1.2 (65.48\%) and accepting 86,412 peptide-spectrum matches at 80\% precision.
\item \textbf{High-Throughput Non-Autoregressive Inference:} DFlowNovo decodes all residue positions in parallel within a fixed budget of $K=20$ integration steps, achieving 227.1 to 255.8 spectra per second on an NVIDIA GPU---a 4.3$\times$ to 4.9$\times$ throughput increase over autoregressive beam search.
\item \textbf{Physical Dissection of Fragmentation Chemistry:} We analyze the biological and mass spectrometry limits of sequencing, proving that 30.88\% of remaining errors on human benchmarks originate from the fundamental isobaric equivalence of Leucine and Isoleucine ($113.084\,\text{Da}$), which cannot be resolved without radical-driven $w$-ion dissociation.
\end{enumerate}

The remainder of this thesis is structured as follows. Chapter 2 details the mathematical foundations of continuous-time discrete flow matching and Continuous-Time Markov Chains. Chapter 3 presents the architectural design of DFlowNovo and the exact dynamic programming knapsack reachability filter. Chapter 4 reports extensive empirical benchmarks, latency comparisons, and ablation experiments. Chapter 5 provides an in-depth physical and chemical error analysis. Finally, Chapter 6 summarizes the findings and outlines promising future directions.
